\newpage
\section{NumPy}

\begin{idea}
A NumPy \code{ndarray} is a grid of values that all have the \textbf{same type} (dtype). Operations apply to the whole array at once (``vectorized''), with no Python loop. That is why NumPy is fast.
\end{idea}

\subsection{Creating arrays and their attributes}
\begin{codebox}
\begin{Verbatim}
import numpy as np
np.array([1, 2, 3])                 # from a list
np.array([[1,2,3],[4,5,6]])         # 2D: 2 rows, 3 columns
np.zeros((2,3)), np.ones((2,3)), np.full((2,2), 7), np.empty(3)
np.identity(3)  /  np.eye(3)        # ONES on the diagonal, ZEROS elsewhere
np.arange(0, 10, 2)                 # [0 2 4 6 8]  evenly spaced, like range()
np.linspace(0, 1, 5)                # 5 evenly spaced points from 0 to 1 (end included)
A = np.array([[1,2,3],[4,5,6]])
A.shape   # (2, 3)   tuple of dimensions
A.size    # 6        TOTAL number of elements (2*3)
A.ndim    # 2        number of dimensions
len(A)    # 2        length of the FIRST dimension only
A.dtype   # int64
A.astype(float)                     # convert the type (returns a new array)
\end{Verbatim}
\end{codebox}
\textbf{dtype rules:}
\begin{itemize}
  \item An array has \textbf{one type only}. Mixed values are upcast: \code{[1, 2.5]} becomes float64, and \code{[1, \textquotesingle{}a\textquotesingle{}]} becomes strings (\code{\textless{}U21}).
  \item With \code{dtype=object} an array can hold anything. In pandas, a DataFrame whose columns have different types gives \code{df.values.dtype == object}.
  \item An explicit dtype is kept: \code{np.array([[[1,2],[3,4],[5,6]]], dtype=np.uint8)} has shape \textbf{(1, 3, 2)} and dtype \textbf{uint8}. (Count the brackets: 1 block, containing 3 rows of 2.)
  \item \code{np.array([1,2,3],[4,5,6],[7,8,9])} is an \textbf{error}: it passes 3 separate lists instead of one list of lists, so the answer is ``None of the options''.
\end{itemize}

\textbf{Random numbers}
\begin{codebox}
\begin{Verbatim}
np.random.seed(42)            # same "random" numbers every run (reproducible)
np.random.rand(4, 3)          # uniform in [0,1), shape (4,3)
np.random.randn(4, 3)         # standard normal N(0,1), shape (4,3)
np.random.randint(0, 10, 5)   # 5 integers in [0, 10)  (10 excluded)
np.random.random(100)         # 100 uniform values, shape (100,)
np.random.uniform(a, b, n); np.random.choice(arr, k); np.random.permutation(arr)
np.random.shuffle(arr)        # shuffles IN PLACE
\end{Verbatim}
\end{codebox}

\subsection{Indexing and slicing}
\begin{codebox}
\begin{Verbatim}
arr = np.array([[10,20,30],
                [40,50,60],
                [70,80,90]])
arr[0, 1]          # 20       (row 0, column 1);  arr[0][1] is the same
arr[1]             # [40 50 60]   a whole row
arr[:, 1]          # [20 50 80]   a whole column
arr[:2, 1:]        # [[20 30],[50 60]]   rows 0-1, columns 1-2
a[:3, -5:]         # first 3 rows, LAST 5 columns
\end{Verbatim}
\end{codebox}

\textbf{Fancy indexing} (indexing with lists of integers):
\begin{codebox}
\begin{Verbatim}
arr[[0, 2], [1, 2]]     # pairs (0,1) and (2,2)  -> [20 90]   1D result!
arr[[0, 2]]             # rows 0 and 2            -> [[10 20 30],[70 80 90]]
arr[[1,0]][:, [2,0]]    # first take rows 1,0 (in THAT order), then columns 2,0
\end{Verbatim}
\end{codebox}
For 23H Q16, \code{arr[[1,4,2]][:, [0,4,2]]} takes rows 1, 4, 2 and then columns 0, 4, 2 \emph{in that order}. The result is a 3$\times$3 block.

\textbf{Boolean indexing}
\begin{codebox}
\begin{Verbatim}
arr[[True, False, True]]   # rows 0 and 2 -> [[10 20 30],[70 80 90]]
arr[[False, True, False]]  # -> array([[40, 50, 60]])   (still 2D)
arr[arr > 50]              # [60 70 80 90]   1D array of matching values
arr[arr > 50] = 0          # sets those elements to 0
np.where(arr > 50, 1, 0)   # element-wise "if": 1 where True, else 0
(arr == 0).sum()  ==  np.count_nonzero(arr == 0)   # counts the True values
\end{Verbatim}
\end{codebox}
Combine conditions with \code{\&} (and), \code{|} (or) and \code{\textasciitilde{}} (not), with brackets around each condition: \code{arr[(arr \textgreater{} 10) \& (arr \textless{} 50)]}.

\begin{idea}
\textbf{Views and copies.} A slice (\code{arr[1:3]}) is a \textbf{view}: changing it changes the original array. Fancy or boolean indexing returns a \textbf{copy}. Use \code{.copy()} when you want an independent array.
\end{idea}

\textbf{Overlapping in-place operation} (22H resit): \code{a = np.array([5,2,2,3]); a[1:] += a[:3]}. NumPy computes the right-hand side first, \code{[5,2,2] + [2,2,3] = [7,4,5]}, and then writes it back. The result is \textbf{[5, 7, 4, 5]}.

\subsection{Universal functions (ufuncs) and aggregation over an axis}
Element-wise functions: \code{np.sqrt}, \code{np.exp}, \code{np.log}, \code{np.abs}, \code{np.maximum(a,b)}, \code{np.add}, \code{np.multiply}, \code{np.modf} (fractional and integer parts), \code{np.logical\_xor}. Arithmetic \code{a + b}, \code{a * b} and \code{a ** 2} is also element-wise. Note that \code{*} is \emph{not} matrix multiplication.

Aggregations: \code{sum}, \code{mean}, \code{std}, \code{var}, \code{min}, \code{max}, \code{argmin}/\code{argmax} (position of the min/max), \code{cumsum}, \code{cumprod}, \code{any}, \code{all}.

\begin{idea}
\textbf{axis = the dimension that disappears.}
\begin{itemize}
  \item \code{axis=0} collapses the rows, so you get \textbf{one value per column}.
  \item \code{axis=1} collapses the columns, so you get \textbf{one value per row}.
\end{itemize}
For shape \code{(2,3,4)}: \code{sum(axis=0)} has shape \code{(3,4)}, \code{sum(axis=1)} has shape \code{(2,4)}, and \code{sum(axis=2)} has shape \code{(2,3)}.
\end{idea}
\begin{codebox}
\begin{Verbatim}
arr = np.array([[1,2,3],[4,5,6],[7,8,9]])
arr.sum()          # 45
arr.sum(axis=0)    # [12 15 18]   column sums
arr.sum(axis=1)    # [ 6 15 24]   row sums          <- 25H Q36
arr.mean(axis=1)   # [2. 5. 8.]
arr.cumsum(axis=0) # running totals down each column
\end{Verbatim}
\end{codebox}

\subsection{Broadcasting}
\begin{idea}
\textbf{Broadcasting} lets arrays of different shapes be combined. Compare the shapes \textbf{from the right}. Two dimensions are compatible when they are \textbf{equal} or one of them is \textbf{1}. A missing dimension counts as 1. If any pair is incompatible, you get a \textbf{ValueError}.
\end{idea}
\begin{codebox}
\begin{Verbatim}
a = np.array([[1,2,3],[4,5,6]])    # shape (2,3)
b = np.array([10,20,30])           # shape (3,)  -> treated as (1,3)
a + b    # [[11 22 33],[14 25 36]]    b is added to EVERY ROW     <- 25H Q27
a + 10   # adds 10 to every element (a scalar always broadcasts)
(3,2) / (2,)      OK : divides each row element-wise
(2,3) + (2,)      ERROR: 3 vs 2 are not compatible
(4,1) + (3,)      OK : the result has shape (4,3)
\end{Verbatim}
\end{codebox}

\subsection{Reshape, transpose, swapaxes}
\begin{codebox}
\begin{Verbatim}
arr = np.arange(12)
arr.reshape(3, 4)     # returns a NEW shape; arr itself is NOT changed
arr = arr.reshape(3, 4)   # you must reassign to keep it
arr.reshape(-1, 2)    # -1 = "compute this dimension for me" -> (6, 2)
arr.ravel() / arr.flatten()   # back to 1D (ravel = view when possible, flatten = copy)
arr.T  /  np.transpose(arr)   # 2D transpose: rows <-> columns
\end{Verbatim}
\end{codebox}
\textbf{transpose with axes:} \code{x.transpose(i, j, k)} puts old axis $i$ first, then $j$, then $k$. The new shape is \code{(shape[i], shape[j], shape[k])}.
\begin{itemize}
  \item Shape \code{(5,2,4)} with \code{.transpose(2,0,1)} gives \code{(4,5,2)} (22H resit Q27).
  \item To swap the first and last axes of \code{(2,3,4)} you can use \code{arr.transpose(2,1,0)} \textbf{and} \code{np.swapaxes(arr,0,2)}, so the answer is ``1\&3'' (25H Q16).
  \item \code{arr.transpose((1,0,2))} is the same as \code{arr.swapaxes(0,1)} (2022 sample).
\end{itemize}
\textbf{22H resit Q34:} \code{arr = np.random.random(100); arr.reshape(10,10)} leaves the shape as \textbf{(100,)}, because the result was not reassigned.

\subsection{Combining and splitting}
\begin{codebox}
\begin{Verbatim}
a, b = np.array([1,2]), np.array([3,4])
np.concatenate((a, b))     # [1 2 3 4]   joins along an EXISTING axis (default axis=0)
np.stack((a, b))           # [[1 2],[3 4]]  adds a NEW axis -> shape (2,2)
np.vstack((a, b)), np.hstack((a, b))   # vertical / horizontal stacking
np.split(arr, 3), np.vsplit, np.hsplit
np.tile(a, 2) -> [1 2 1 2];   np.repeat(a, 2) -> [1 1 2 2]
\end{Verbatim}
\end{codebox}
For 2D arrays, \code{np.concatenate([x, y])} with the default \code{axis=0} stacks them \textbf{vertically}: it adds rows. With \code{axis=1} it adds columns.

\subsection{Sorting, searching, set logic}
\begin{codebox}
\begin{Verbatim}
arr = np.array([-1, 2, 4, 0, 2, -1, 1, 0])
np.sort(arr)       # returns a SORTED COPY (ascending); arr is UNCHANGED
arr.sort()         # sorts arr IN PLACE
np.sort(arr)[::-1] # descending
arr.argsort()      # the INDICES that would sort the array
data[data[:, 1].argsort()]   # sort the rows of a 2D array by column 1
np.partition(arr, 3)  # elements smaller than the 3rd smallest go to the left (not fully sorted)
np.searchsorted(sorted_arr, v)   # where v would be inserted
np.unique(arr)        # sorted unique values
np.isin(arr, [0, 1])  # element-wise "is in" -> boolean array
np.intersect1d, np.union1d, np.setdiff1d
\end{Verbatim}
\end{codebox}
\textbf{22H resit Q40:} after \code{np.sort(arr)} with no assignment, arr is stored in its \textbf{original order}.

\subsection{Linear algebra and file I/O}
\begin{codebox}
\begin{Verbatim}
a.dot(b) == a @ b == np.dot(a, b)    # matrix / dot product: ALL equivalent
a * b                                # element-wise (NOT the matrix product)
np.linalg.inv(A), np.linalg.det(A)
np.save("x.npy", arr);  np.load("x.npy")         # binary .npy file
np.savez("x.npz", a=arr1, b=arr2)                # several arrays in one file
np.savez_compressed(...); np.savetxt / np.loadtxt  # text files
\end{Verbatim}
\end{codebox}
To \emph{write} an array to a file you use \code{np.save}. If the options are only \code{read\_csv}, \code{load} and \code{open}, the answer is ``None of the options''.

\begin{traps}
\begin{itemize}
  \item \code{[20, 90]} and \code{[[20, 90]]} are different: pairs of index lists give a 1D result.
  \item \code{.reshape()} and \code{np.sort()} do not change the array unless you reassign the result.
  \item \code{np.identity} is \emph{not} random, and it is \emph{not} all ones.
  \item \code{.size} counts the elements, while \code{.shape} is a tuple.
  \item \code{arr.transform(...)}, \code{arr.T(0,1)} and \code{arr.reshape(0,1,2)} are not ways to swap axes.
\end{itemize}
\end{traps}
